% STATUS: written 2026-11-06 (English gate clean; sources read in full as recorded below).
% Sources read in full for the cited passages: Biswas-Gurjar-Kolte 2014 (pp. 270--273), Graf 2022 (pp. 115--117), Granger-Schulze 2014 (pp. 1--2).
% Sources: BiswasGurjarKolte2014, Graf2022LZPosChar, GrangerSchulze2014Negative.
\section{The Zariski--Lipman conjecture}\label{sec:commutative:zariski-lipman}

Each derivation below is of a commutative algebra, the setting the Zariski-Lipman conjecture is stated in.

Section~\ref{sec:commutative:kahler} identified the derivations of a
commutative algebra with the dual of one module. Lipman asked in 1965 how much
this module says about the singularities of the algebra
\cite[p.~270]{BiswasGurjarKolte2014}. The paper is
\emph{Free derivation modules on algebraic varieties}
\cite[p.~270]{BiswasGurjarKolte2014}.

The question is this~\cite[p.~270]{BiswasGurjarKolte2014}. Suppose that the
module of derivations of a local ring is free. Must the ring then come from a
smooth point~\cite[p.~270]{BiswasGurjarKolte2014}? The question is now named
after Zariski and Lipman \cite[p.~270]{BiswasGurjarKolte2014}.

\begin{theorem}[{\cite[Conjecture~1.1, p.~270]{BiswasGurjarKolte2014}}]\label{thm:comm:zl-statement}
Let $V$ be an algebraic variety over a field $k$ of characteristic $0$, let
$p$ be a closed point of $V$ and let $R$ be the local ring of $V$ at $p$. If
the module of $k$-derivations $\operatorname{Der}_k(R)$ is a free $R$-module,
then $V$ is smooth at $p$.
\end{theorem}

The statement of Theorem~\ref{thm:comm:zl-statement} is a
conjecture~\cite[p.~270]{BiswasGurjarKolte2014}, quoted as it stands in the
literature, and no proof is claimed here. Its hypotheses carry
weight~\cite[p.~270]{BiswasGurjarKolte2014}. Lipman proved that if
$\operatorname{Der}_k(R)$ is $R$-free, then $R$ is
normal~\cite[p.~270]{BiswasGurjarKolte2014}. So the question concerns only
normal varieties. If the ring is not normal, the answer can be negative
already in dimension one~\cite[p.~270]{BiswasGurjarKolte2014}.

Several cases are settled~\cite[p.~270]{BiswasGurjarKolte2014}. Biswas, Gurjar
and Kolte collect the main
ones~\cite[p.~270]{BiswasGurjarKolte2014}. Hochster proved the conjecture for
positively graded domains, and Scheja and Storch proved it for
hypersurfaces~\cite[p.~270]{BiswasGurjarKolte2014}. Becker reduced it to the
case of isolated singularities~\cite[p.~270]{BiswasGurjarKolte2014}, and van
Straten and Steenbrink settled that case in dimension at least
three~\cite[p.~270]{BiswasGurjarKolte2014}.

Flenner settled non-isolated singularities of dimension at least three in
codimension at least three~\cite[p.~270]{BiswasGurjarKolte2014}. Källström
proved the conjecture for complete
intersections~\cite[p.~270]{BiswasGurjarKolte2014}. Druel and, independently,
Graf proved it for local rings with log canonical
singularities~\cite[p.~270]{BiswasGurjarKolte2014}.

Part of this list is recalled in the cited paper as a single
statement~\cite[Theorem~2.2, pp.~272--273]{BiswasGurjarKolte2014}. That
statement is Theorem~\ref{thm:comm:zl-cases} below. It is the form later work
quotes, because it gathers the settled cases into one hypothesis.

\begin{theorem}[{\cite[Theorem~2.2, pp.~272--273]{BiswasGurjarKolte2014}}]\label{thm:comm:zl-cases}
Let $V$ be an affine algebraic variety such that the algebraic tangent
bundle of the smooth locus $V^0$ is trivial. Then $V$ is smooth in each of
the following cases: the coordinate ring of $V$ is a positively graded
domain; $V$ is a complete-intersection singularity; or $V$ is a surface with
at most rational singularities.
\end{theorem}

For surfaces the conjecture is nearly closed. Biswas, Gurjar and Kolte prove
it for a wide class of affine surfaces: those whose smooth locus has
logarithmic Kodaira dimension at most one
\cite[Theorem~1.2, p.~270]{BiswasGurjarKolte2014}.

Three authors worked on the rest: Lipman, Becker and
Flenner~\cite[p.~115]{Graf2022LZPosChar}. Graf reports their work, and says
that the case of a normal surface is essentially the only open case left
\cite[p.~115]{Graf2022LZPosChar}.

In prime characteristic the picture
changes~\cite[p.~115]{Graf2022LZPosChar}. Lipman himself found counterexamples,
and Graf recalls them~\cite[p.~116]{Graf2022LZPosChar}.

\begin{example}[{\cite[p.~116]{Graf2022LZPosChar}}]\label{ex:comm:lipman}
Let $k$ be a perfect field of characteristic $p$, and let
$X\subset\mathbb{A}^3_k$ be the surface defined by $xy - z^n = 0$, where $p$
divides $n$. Then the tangent sheaf of $X$ is freely generated by the two
derivations $\delta_1 = x\partial_x - y\partial_y$ and
$\delta_2 = \partial_z$, although $X$ is singular.
\end{example}

So Graf studies in the cited paper how far the failure in prime characteristic
goes \cite[p.~115]{Graf2022LZPosChar}. He proves that for rational double
points the conjecture holds except for a short finite list of cases. Among
rational double points in characteristic $p \ge 7$, he adds, the counterexample
of Lipman is the only one \cite[p.~115]{Graf2022LZPosChar}.

A second question about the derivations of a singularity is qualitative.
Give a quasihomogeneous singularity its positive grading. Does some grading
exist for which no derivation lowers the degree
\cite[p.~1]{GrangerSchulze2014Negative}? Wahl made the conjecture that it
does, for every quasihomogeneous isolated normal singularity
\cite[Conjecture, p.~2]{GrangerSchulze2014Negative}. Granger and Schulze
settle the conjecture for isolated complete intersections of order at least
three (Theorem~\ref{thm:comm:wahl-order}) \cite[Theorem~1.1, p.~2]{GrangerSchulze2014Negative}.

\begin{theorem}[{\cite[Theorem~1.1, p.~2]{GrangerSchulze2014Negative}}]\label{thm:comm:wahl-order}
Let $A$ be a quasihomogeneous normal isolated complete intersection
singularity of order at least three. Then $A$ has no derivation of negative
weighted degree with respect to any positive grading.
\end{theorem}

They also settle every isolated complete intersection of embedding dimension
at most five \cite[Theorem~1.3, p.~2]{GrangerSchulze2014Negative}. Above
embedding dimension five the conjecture fails, and they give a family of
counterexamples \cite[Example~1.2, pp.~1--2]{GrangerSchulze2014Negative}. The
first has embedding dimension six, and on it the derivation $\eta$ of degree
$-1$ is written out explicitly \cite[Example~1.2, p.~2]{GrangerSchulze2014Negative}.
So the sign of the degree of a derivation separates the cases, and the
separation is exact at embedding dimension five
\cite[Theorem~1.3, p.~2]{GrangerSchulze2014Negative}.

So Theorem~\ref{thm:comm:zl-statement} compares an algebraic property with a
geometric one~\cite[p.~270]{BiswasGurjarKolte2014}: the module of derivations
is free, and the variety is smooth.

The conjecture of Theorem~\ref{thm:comm:zl-statement} is open, and this
section records it as the first problem of the chapter. % claim-ok: own signpost introducing the moved problem, which carries its own citation

\begin{problem}[{\cite[Conjecture~1.1, p.~270]{BiswasGurjarKolte2014}}]\label{prob:open:zl}
Let $V$ be an algebraic variety over a field $k$ of characteristic $0$.
Let $p$ be a closed point of $V$, and let $R$ be the local ring of $V$ at $p$.
If the module of $k$-derivations $\Der_k(R)$ is a free $R$-module, then $V$ is
smooth at $p$.
\end{problem}

The conjecture is settled for positively graded domains, for hypersurfaces, for
complete intersections and for log canonical singularities
\cite[p.~270]{BiswasGurjarKolte2014}. Biswas, Gurjar and Kolte settle it for
every affine surface whose smooth locus has logarithmic Kodaira dimension at
most one \cite[Theorem~1.2, p.~270]{BiswasGurjarKolte2014}. They leave one
case of the projective problem open. % claim-ok: own signpost introducing the second moved problem, whose open case the cited theorem states

\begin{problem}[{\cite[Remark~1.3, p.~271]{BiswasGurjarKolte2014}}]\label{prob:open:zl-projective}
Let $V$ be a projective surface whose smooth locus has a trivial tangent bundle
and logarithmic Kodaira dimension $-\infty$.
Suppose the base of the $\mathbb{P}^1$-fibration of a resolution of $V$ is an
elliptic curve, and the exceptional divisor has only reduced feathers.
Is $V$ then smooth?
\end{problem}

Problem~\ref{prob:open:zl} is a statement about characteristic $0$ only.
The conjecture fails in prime characteristic
\cite[p.~116]{Graf2022LZPosChar}, and the example of Lipman above is the
smallest failure recorded in this section. % claim-ok: own signpost from the moved problem to the counterexample above

The next section keeps the algebraic side of the conjecture, and asks for
derivations that are locally nilpotent
\cite[Theorem, p.~47]{FerreroLequainNowicki1992}. From those it reads
polynomial coordinates off.
% claim-ok: own signpost to the next section, moved after the problems
